% Источники: exam/materials/моделирование.pdf, разделы 25--28; https://github.com/HumsterProgrammer/iu9-notes/blob/main/8sem/Моделирование/Моделирование_лекция-04_19-02-26.md.
\section{Организация квазипараллелизма в рамках последовательной архитектуры компьютера.}

\textbf{Управляющая программа моделирования} -- это компонент имитационной
модели, который координирует выполнение всех процессов, обеспечивает
синхронизацию, управляет временем моделирования и собирает статистику.

Функции, реализуемые УПМ в специализированных средах моделирования:
\begin{enumerate}
  \item внутренняя синхронизация компонент модели;
  \item внешняя синхронизация компонент модели;
  \item уточнение исходной информации при моделировании;
  \item контроль за ходом имитации;
  \item организация окончания имитации;
  \item организация сбора статистики;
  \item документирование.
\end{enumerate}

Приставка \textbf{квази} отражает последовательный характер обслуживания
событий в модели, одновременно возникающих в разных компонентах реальной
системы.

Различают пять способов организации квазипараллелизма:
\begin{enumerate}
  \item имитация активностями;
  \item событийный подход;
  \item транзактный подход;
  \item процессный подход;
  \item агрегатный подход.
\end{enumerate}

\subsection*{Подход на основе активностей}

Имитация на основе активностей -- это метод моделирования, при котором
система описывается через набор активностей, выполняемых ее компонентами.
Каждая активность представляет собой действие, которое занимает определенное
время и приводит к изменению состояния системы.

Вводится понятие активности
\[
  A_{ij} = \bigl\{(Al_{ij}, \tau_{ij})\bigr\},
\]
представляющее собой кортеж, элементами которого являются алгоритм
$j$-ого события для $i$-ого объекта и момент модельного времени. Вместо
$\tau_{ij}$ можно рассматривать изменение времени с предыдущего момента.

Достоинства: вычислительно простая схема, учитывающая взаимодействие
объектов; высокая точность моделирования. Недостаток: если количество
объектов большое, решение вычислительной задачи для каждого объекта
приводит к резкому росту объема вычислений.

Пример применения: система с небольшим количеством компонент и состояний,
где для каждого состояния известен свой алгоритм обработки.

\subsection*{Подход на основе расписания событий}

Событие представляет собой мгновенное изменение состояния некоторого
объекта системы, то есть изменение значений его атрибутов. Окончание любой
активности в системе является событием, так как приводит к изменению
состояния объекта и может служить инициатором другой работы в системе.

Когда объектов очень много, $\Delta\tau_{ij}$ различны, алгоритмом
функционирования объекта пренебрегают. Реализация модели основывается на
расписании. Для реализации модели необходимо построить запрос, результатом
которого является плоская таблица расписания событий.

Пример применения: расписание занятий, когда для решения задачи достаточно
знать, где, когда и какое занятие состоялось, без подробного описания
заполнения аудитории.

\subsection*{Агрегатный подход}

Если в системе имеет место тесное взаимодействие компонент, то компоненты
обмениваются между собой сигналами. Каждый выходной сигнал от одной
компоненты является входным сигналом для другой компоненты.

Если функциональные действия аппроксимируются явно задаваемыми
математическими зависимостями, позволяющими определить момент появления
выходных сигналов при наличии входных сигналов, поступающих от других
компонентов, то создаются условия для построения имитационной модели по
модульному принципу. В этом случае каждый из модулей имитационной модели
строится по унифицированной структуре и называется \textbf{агрегатом}.

Агрегат является математической схемой, с помощью которой возможно описание
большого круга реальных процессов в системе. В любой момент $t_{ij}$ агрегат
может находиться в одном из возможных состояний. Состояния агрегата $Z$
являются функциями времени. Переход агрегата $Z$ из состояния в состояние
описывается оператором перехода $H$, который позволяет по предыдущему
состоянию определить очередное его состояние.

Агрегат имеет входы, куда поступают входные сигналы $X(t)$ от других
агрегатов, и выходы, на которых формируются выходные сигналы $Y(t)$. Кроме
того, у агрегата имеются дополнительные входы, на которые поступают
управляющие сигналы $g(t)$. Выходные сигналы $Y(t)$ формируются из входных
$X(t)$ и управляющих $g(t)$ сигналов оператором выхода $G$ в результате его
взаимодействия с оператором $H$.

\clearpage
